% Generated draft: requires mathematical and source review.
\documentclass[11pt]{article}
\usepackage[margin=25mm]{geometry}
\usepackage{fontspec}
\setmainfont{Latin Modern Roman}
\usepackage{amsmath,amssymb,mathtools,unicode-math}
\setmathfont{Latin Modern Math}
\usepackage{enumitem,needspace,xurl,hyperref}
\setlength{\parindent}{0pt}
\setlength{\parskip}{5pt}
\setlength{\emergencystretch}{3em}
\newcommand{\N}{\mathbb{N}}
\newcommand{\Z}{\mathbb{Z}}
\newcommand{\Q}{\mathbb{Q}}
\newcommand{\R}{\mathbb{R}}
\newcommand{\C}{\mathbb{C}}
\newcommand{\F}{\mathbb{F}}
\newcommand{\A}{\mathbb{A}}
\newcommand{\PP}{\mathbb{P}}
\newcommand{\E}{\mathbb{E}}
\newcommand{\eps}{\varepsilon}
\newcommand{\dd}{\,\mathrm{d}}
\begin{document}
\section*{UnsolvedMath ID: 701176}
\textbf{Source problem code:} catalogue entry
\subsection{Definitions and mathematical statement}
Let \(\mathfrak A\) be the variety of abelian groups and \(\mathfrak N_c\) the variety of nilpotent groups of class at most \(c\). Under the standard convention for products of group varieties, \(G\in\mathfrak A\mathfrak N_c\) means that there exists an abelian normal subgroup \(A\trianglelefteq G\) such that \(G/A\in\mathfrak N_c\).

Let \(F_r=F_r(\mathfrak A\mathfrak N_c)\) be the relatively free group of rank \(r\) in this variety. Thus a free generating set \(X\) of cardinality \(r\) has the property that every map from \(X\) to a group in \(\mathfrak A\mathfrak N_c\) extends uniquely to a homomorphism from \(F_r\). An element \(g\in F_r\) is \emph{primitive} if it belongs to some free basis of \(F_r\).

\textbf{Question.} Does there exist an algorithm recognizing primitive elements of \(F_r\)? Under the usual interpretation with fixed positive integers \(c\) and \(r\), fix a free basis \(X=\{x_1,\ldots,x_r\}\). The requested algorithm must halt on every finite word \(w\) in \(X^{\pm1}\) and decide whether its value \(g=w(X)\in F_r\) satisfies
\[
\exists\,(b_1,\ldots,b_r)\text{ a free basis of }F_r\quad b_1=g.
\]
The supplied source does not explicitly specify the ranges of \(c\) and \(r\), whether infinite rank or the boundary case \(c=0\) is included, or whether a single algorithm uniform in \(c\) and \(r\) is required. The finite-rank, fixed-parameter formulation above is therefore an explicit interpretation requiring review.
\subsection{Short English statement}
The problem asks whether one can decide, from a word in a relatively free abelian-by-nilpotent group, whether the represented element belongs to a free basis. The intended parameter ranges and any requirement of uniformity remain unspecified.
\subsection{Sources}
\begin{itemize}
\item Local UnsolvedMath registry, permanent ID \(701176\), exact source problem code \texttt{catalogue entry}, titled \emph{Kourovka Notebook Problem 15.88}. The supplied record labels the problem open; this status has not been independently verified.
\item Attributed in the supplied record to E. I. Timoshenko, proposed in 2002. The stated source is \emph{The Kourovka Notebook (Unsolved Problems in Group Theory)}, 15th issue (2002), Problem 15.88, as carried in the supplied 21st-issue edition reference (2026), arXiv:1401.0300v48. The collection description names E. I. Khukhro and V. D. Mazurov as editors. These bibliographic details are supplied attribution and have not been independently checked.
\end{itemize}
\end{document}
